\ifdefined\gkver\else\def\gkver{student}\fi
\documentclass[\gkver]{gaokaozhenti}
\begin{document}

\chapter{2025年海南}

\begin{enumerate}
\item
%% number: 1
%% paperName: 2025年海南高考第 1 题 3 分
%% typeId: 1
%% type: 单选题
%% chapter: 原子和原子核
%% point: 天然放射性
%% method: 无
%% score: 3
%% degree: 950
%% duplicateId: 0
%% body:
海南昌江核电示范项目“玲珑一号”是中国自主研发的全球首个陆上商用模块化小型核反应堆，其中发生的一个核反应方程为 \ce{^{235}_{92} U + X -> ^{144}_{56} Ba + ^{89}_{36} Kr + 3^{1}_{0} n} 中，则 X 是\xzanswer{C}

\fourchoices[answer=C]
{\ce{^{4}_{2} He}}
{\ce{^{1}_{1} H}}
{\ce{^{1}_{0} n}}
{\ce{^{0}_{-1} e}}

%% answer: C
%% memo:
\memoanswer{
	根据质量数守恒和核电荷数守恒可知 X 是 \ce{^{1}_{0} n}。

	故选 C。
}

\item
%% number: 2
%% paperName: 2025年海南高考第 2 题 3 分
%% typeId: 1
%% type: 单选题
%% chapter: 电场
%% point: 静电感应
%% method: 无
%% score: 3
%% degree: 850
%% duplicateId: 0
%% body:
带正电的金属球靠近不带电验电器金属小球 a，则关于验电器金属小球 a 和金属箔 b，下列说法正确的是\xzanswer{C}

\onepicture[w=4.30cm]{figs/02.png}

\fourchoices[answer=C]
{a、b 都带正电}
{a、b 都带负电}
{a 带负电、b 带正电}
{a 带正电、b 带负电}

%% answer: C
%% memo:
\memoanswer{
	由图可知，验电器本来不带电，由于同种电荷相互排斥，异种电荷相互吸引，正金属球靠近不带电验电器金属小球 a，使得金属球 a 带负电荷，从而导致金属箔 b 带上正电荷。

	故选 C。
}

\item
%% number: 3
%% paperName: 2025年海南高考第 3 题 3 分
%% typeId: 1
%% type: 单选题
%% chapter: 直线运动
%% point: 运动图象
%% method: 无
%% score: 3
%% degree: 750
%% duplicateId: 0
%% body:
ETC 是电子不停车收费系统的简称，常见于高速公路出入口，只要在车挡风玻璃上安装一个打卡装置，就能实现快速收费，提高通行效率。如图所示是一辆汽车通过 ETC 通道运动过程的速度-时间图像，其中 $t_{1} \sim t_{2}$ 时间内的图线是一条平行于 $t$ 轴的直线，则\xzanswer{A}

\onepicture[w=5.40cm]{figs/03.png}

\fourchoices[answer=A]
{汽车在 $0 \sim t_{1}$ 时间内做匀减速直线运动}
{汽车在 $t_{1} \sim t_{2}$ 时间内处于静止状态}
{汽车在 $0 \sim t_{1}$ 和 $t_{2} \sim t_{3}$ 时间内的加速度方向相同}
{汽车在 $0 \sim t_{1}$ 和 $t_{2} \sim t_{3}$ 时间内的速度方向相反}

%% answer: A
%% memo:
\memoanswer{
	A．由图可知 $v-t$ 图像的斜率表示加速度，$0 \sim t_{1}$ 时间内加速度为负且恒定，速度为正，加速度方向与速度方向相反，故 $0 \sim t_{1}$ 时，汽车做匀减速直线运动，故 A 正确；

	B．$t_{1} \sim t_{2}$ 内，汽车做匀速直线运动，故 B 错误；

	C．$0 \sim t_{1}$ 内加速度为负，$t_{2} \sim t_{3}$ 内加速度为正，故 $0 \sim t_{1}$ 和 $t_{2} \sim t_{3}$ 内，汽车加速度方向相反，故 C 错误；

	D．$0 \sim t_{1}$ 和 $t_{2} \sim t_{3}$ 内，汽车速度方向相同，均为正，故 D 错误。

	故选 A。
}

\item
%% number: 4
%% paperName: 2025年海南高考第 4 题 3 分
%% typeId: 1
%% type: 单选题
%% chapter: 曲线运动
%% point: 人造地球卫星
%% method: 无
%% score: 3
%% degree: 650
%% duplicateId: 0
%% body:
2025 年 4 月 24 日，载人飞船神州二十号在酒泉卫星发射中心点火发射，在进入预定轨道后成功与空间站完成自主快速交会对接，然后绕地球做匀速圆周运动。已知空间站轨道高度低于地球同步卫星轨道，则下面说法正确的是\xzanswer{B}

\fourchoices[answer=B]
{火箭在加速升空的过程中处于失重状态}
{航天员在空间站所受地球的引力小于其在地面上受到的地球引力}
{空间站绕地球做匀速圆周运动的角速度小于地球自转角速度}
{空间站绕地球做匀速圆周运动的加速度小于地球同步卫星的加速度}

%% answer: B
%% memo:
\memoanswer{
	A．火箭加速升空过程，加速度方向竖直向上，则处于超重状态，故 A 错误；

	B．根据 $F = \frac{GMm}{R^{2}}$，宇航员与地球的质量不变，宇航员在空间站离地心更远，则受到的万有引力小于在地表受到万有引力，故 B 正确；

	C．根据 $\frac{GMm}{R^{2}} = m\omega^{2}R$ 可得 $\omega = \sqrt{\frac{GM}{R^{3}}}$，可知空间站绕地球做匀速圆周运动的角速度大于同步卫星的角速度，即大于地球自转角速度，故 C 错误；

	D．根据 $\frac{GMm}{R^{2}} = ma$ 可得 $a = \frac{GM}{R^{2}}$，可知空间站绕地球做匀速圆周运动的加速度大于地球同步卫星的加速度，故 D 错误。

	故选 B。
}

\item
%% number: 5
%% paperName: 2025年海南高考第 5 题 3 分
%% typeId: 1
%% type: 单选题
%% chapter: 机械波
%% point: 波的图象
%% method: 无
%% score: 3
%% degree: 550
%% duplicateId: 0
%% body:
如图所示，实线和虚线分别是沿着 $x$ 轴正方向传播的一列简谐横波在 $t = 0$ 时刻和 $t = 0.5 \Us$ 的波形图，已知波的周期 $T > 0.5 \Us$，则下列关于该列波说法正确的是\xzanswer{D}

\onepicture[w=6.20cm]{figs/05.png}

\fourchoices[answer=D]
{波长为 $5 \Ucm$}
{波速为 $10 \Ucms$}
{周期为 $1 \Us$}
{$t$ = 0 时刻，质点 M 向下振动}

%% answer: D
%% memo:
\memoanswer{
	A．由图可知波长为 $10 \Ucm$，故 A 错误；

	BC．$t = 0$ 时刻到 $t = 0.5 \Us$ 的过程中，

	$\Delta t = nT + \frac{3}{4}T,\quad n = 0,1,2\ldots$

	已知波的周期 $T > 0.5 \Us$，则 $n$ 只能取 0，则 $T = \frac{2}{3} \Us$

	波速为 $v = \frac{\lambda}{T} = 15 \Ucms$，故 BC 错误；

	D．简谐横波沿着 $x$ 轴正方向传播，根据同侧法可知 $t = 0$ 时刻，质点 M 向下振动，故 D 正确。

	故选 D。
}

\item
%% number: 6
%% paperName: 2025年海南高考第 6 题 3 分
%% typeId: 1
%% type: 单选题
%% chapter: 电场
%% point: 电场中的图象
%% method: 无
%% score: 3
%% degree: 450
%% duplicateId: 0
%% body:
某静电场电势 $\varphi$ 在 $x$ 轴上分布如图所示，图线关于 $\varphi$ 轴对称，M、P、N 是 $x$ 轴上的三点，OM = ON；有一电子从 M 点静止释放，仅受 $x$ 方向的电场力作用，则下列说法正确的是\xzanswer{D}

\onepicture[w=5.20cm]{figs/06.png}

\fourchoices[answer=D]
{P 点电场强度方向沿 $x$ 负方向}
{M 点的电场强度小于 N 点的电场强度}
{电子在 P 点的动能小于在 N 点的动能}
{电子在 M 点的电势能大于在 P 点的电势能}

%% answer: D
%% memo:
\memoanswer{
	A．由图可知在 $x$ 正半轴沿 $+x$ 方向电势降低，则电场强度方向沿 $x$ 正方向，故 A 错误；

	B．$\varphi-x$ 图像斜率表示电场强度，由图可知 M 点的电场强度大小等于 N 点的电场强度，方向相反，故 B 错误；

	C．电子在电势低处电势能大，故电子在 P 点的电势能小于在 N 点的电势能，根据能量守恒可知，电子在 P 点的动能大于在 N 点的动能，故 C 错误；

	D．电子在电势低处电势能大，故电子在 M 点的电势能大于在 P 点的电势能，故 D 正确。

	故选 D。
}

\item
%% number: 7
%% paperName: 2025年海南高考第 7 题 3 分
%% typeId: 1
%% type: 单选题
%% chapter: 电路
%% point: 传感器
%% method: 无
%% score: 3
%% degree: 550
%% duplicateId: 0
%% body:
某同学设计的光控照明电路如图所示，直流电源电动势 $E = 3 \UV$（内阻不计），$R_{1}$ 和 $R_{2}$ 其中一个是定值电阻，一个是光敏电阻。光敏电阻的阻值随照度变化如下表，物理学中用照度描述光的强弱，光越强照度越大。当照度低至 $1.0 \Ulx$ 时，1、2 两端电压升至 $2 \UV$，控制开关自动启动照明设备。不考虑控制开关对设计电路的影响，下列判断正确的是\xzanswer{A}

\begin{center}
\begin{tabular}{|c|c|c|c|c|c|}
\hline
照度/$\Ulx$ & 0.6 & 0.8 & 1.0 & 1.2 & $\ldots$ \\
\hline
电阻/$\UkO$ & 28 & 23 & 20 & 18 & $\ldots$ \\
\hline
\end{tabular}
\end{center}

\onepicture[w=5.80cm]{\includesvg{figs/07}}

\fourchoices[answer=A]
{$R_{1}$ 是光敏电阻，$R_{2}$ 是 $10 \UkO$ 的定值电阻}
{$R_{1}$ 是光敏电阻，$R_{2}$ 是 $40 \UkO$ 的定值电阻}
{$R_{2}$ 是光敏电阻，$R_{1}$ 是 $10 \UkO$ 的定值电阻}
{$R_{2}$ 是光敏电阻，$R_{1}$ 是 $40 \UkO$ 的定值电阻}

%% answer: A
%% memo:
\memoanswer{
	若 $R_{1}$ 是光敏电阻，随着光强的减弱，光敏电阻阻值增大，1、2 两端电压不断升高，当照度低至 $1.0 \Ulx$ 时，$R_{1} = 20 \UkO$，1、2 两端电压升至 $2 \UV$，可得：$R_{1}:R_{2} = U_{1}:U_{2} = 2 \UV:(3-2) \UV$，解得 $R_{2} = 10 \UkO$，A 选项正确、B 选项错误；

	若 $R_{2}$ 是光敏电阻，随着光强的减弱，光敏电阻阻值增大，1、2 两端电压不断减小，不会出现 1、2 两端电压不断升高的现象，C、D 选项错误。
}

\item
%% number: 8
%% paperName: 2025年海南高考第 8 题 3 分
%% typeId: 1
%% type: 单选题
%% chapter: 力物体的平衡
%% point: 物体系平衡
%% method: 整体与隔离
%% score: 3
%% degree: 450
%% duplicateId: 0
%% body:
如图所示，光滑圆弧竖直固定，用轻绳连接两个小球 P、Q，两小球套在圆环上且均处于平衡状态，两小球与圆弧的圆心连线夹角分别为 $30^\circ$ 和 $60^\circ$，则两球质量之比 $m_{P}:m_{Q}$ 为\xzanswer{B}

\onepicture[w=4.40cm]{\includesvg{figs/08}}

\fourchoices[answer=B]
{$1:2$}
{$1:\sqrt{3}$}
{$\sqrt{3}:1$}
{$2:1$}

%% answer: B
%% memo:
\memoanswer{
	方法一：以小球 P 和 Q 为系统，根据力矩平衡有

	$m_{\mathrm{P}} g R \cos 30^\circ = m_{\mathrm{Q}} g R \cos 60^\circ$

	可得 $m_{\mathrm{P}}:m_{\mathrm{Q}} = 1:\sqrt{3}$

	方法二：对 P 受力分析，受重力 $m_{P}g$、圆弧轨道的支持力 $N$ 和轻绳的拉力 $T$，将重力和绳子拉力沿着垂直半径方向分解，根据平衡条件可得 $m_{\mathrm{P}} g \cos 30^\circ = T \cos 45^\circ$

	可得 $m_{\mathrm{P}} = \frac{T \cos 45^\circ}{g \cos 30^\circ}$，同理可得 $m_{\mathrm{Q}} = \frac{T \cos 45^\circ}{g \cos 60^\circ}$

	则 $m_{\mathrm{P}}:m_{\mathrm{Q}} = 1:\sqrt{3}$

	故选 B。
}

\item
%% number: 9
%% paperName: 2025年海南高考第 9 题 4 分
%% typeId: 2
%% type: 多选题
%% chapter: 光学
%% point: 光的干涉
%% method: 无
%% score: 4
%% degree: 750
%% duplicateId: 0
%% body:
下列关于光现象描述正确的是\xzanswer{BC}

\fourchoices[answer=BC]
{海市蜃楼是光的衍射现象}
{光电效应表明光具有粒子性}
{水面上的油膜在阳光下呈现彩色条纹是光的干涉现象}
{斜插在水杯中的筷子看起来像被折断是光的反射现象}

%% answer: BC
%% memo:
\memoanswer{
	A．海市蜃楼是地球上物体反射的光经大气折射而形成的虚像。由于不同高度的空气密度不同，导致光在不均匀的大气中传播时发生折射，从而使人看到远处物体的虚像。所以蜃楼是光的折射现象，A 错误；

	B．光电效应是指在光的照射下，某些物质内部的电子会被光子激发出来而形成电流的现象，它表明光具有粒子性，B 正确；

	C．水面面上的油膜在阳光下呈现彩色条纹，是因为油膜的上下表面反射的光相互叠加，形成了干涉现象。不同颜色的光波长不同，干涉条纹的间距也不同，所以会呈现出彩色条纹。因此，水面面上的油膜在阳光下呈现彩色条纹是光的干涉现象，C 正确；

	D．斜插在水杯中的筷子看起来像被折断是光的折射现象，D 错误。

	故选 BC。
}

\item
%% number: 10
%% paperName: 2025年海南高考第 10 题 4 分
%% typeId: 2
%% type: 多选题
%% chapter: 交流电
%% point: 变压器
%% method: 无
%% score: 4
%% degree: 650
%% duplicateId: 0
%% body:
如图所示，理想变压器接在 $u = 220\sqrt{2}\sin 100\pi t$（$\UV$）交流电上，此时电压表读数为 $10 \UV$，电流表读数为 $4 \UA$，电表皆为理想电表，则\xzanswer{AD}

\onepicture[w=5.80cm]{\includesvg{figs/10}}

\fourchoices[answer=AD]
{输入电压有效值为 $220 \UV$}
{原线圈输入电流为 $5.5 \UA$}
{滑动变阻器向下滑动，电压表读数减小}
{滑动变阻器向上滑动，电流表读数增大}

%% answer: AD
%% memo:
\memoanswer{
	A．根据题干信息可知输入电压的最大值 $U_{m} = 220\sqrt{2} \UV$

	输入电压有效值为 $U_{1} = \frac{U_{m}}{\sqrt{2}} = 220 \UV$，故 A 正确；

	B．理想变压器原副线圈电压比 $\frac{U_{1}}{U_{2}} = \frac{n_{1}}{n_{2}}$

	电流比 $\frac{I_{1}}{I_{2}} = \frac{n_{2}}{n_{1}}$

	其中 $U_{2} = 10 \UV$，$I_{2} = 4 \UA$

	可得 $I_{1} = \frac{2}{11} \UA$，故 B 错误；

	C．原线圈输入电压不变，匝数比不变，根据 $\frac{U_{1}}{U_{2}} = \frac{n_{1}}{n_{2}}$，副线圈两端的电压不变，即电压表读数不变，C 错误；

	D．滑动变阻器向上滑动，滑动变阻器接入的阻值减小，副线圈电压不变，根据 $I_{2} = \frac{U_{2}}{R}$ 可知电流增大，即电流表读数增大，D 正确。

	故选 AD。
}

\item
%% number: 11
%% paperName: 2025年海南高考第 11 题 4 分
%% typeId: 2
%% type: 多选题
%% chapter: 功和能
%% point: 功率
%% method: 无
%% score: 4
%% degree: 550
%% duplicateId: 0
%% body:
一起重机将质量为 $m$ 的集装箱由静止匀加速竖直向上提升，加速度为 $a$，重力加速度为 $g$，不计空气阻力，匀加速时间为 $t$，则\xzanswer{AC}

\fourchoices[answer=AC]
{匀加速的最大速度为 $at$}
{集装箱的机械能增加 $\frac{1}{2}mgat^{2}$}
{起重机的最大输出功率为 $m(g + a)at$}
{起重机对集装箱的作用力为 $ma$}

%% answer: AC
%% memo:
\memoanswer{
	A．匀加速的最大速度 $v_{m} = at$，A 正确；

	B．集装箱的动能增加量为 $\Delta E_{k} = \frac{1}{2}mv_{m}^{2} = \frac{1}{2}ma^{2}t^{2}$

	集装箱上升的高度 $h = \frac{1}{2}at^{2}$

	重力势能的增加量为 $\Delta E_{p} = mgh = \frac{1}{2}mgat^{2}$

	集装箱的机械能增加 $\Delta E = \Delta E_{k} + \Delta E_{p} = \frac{1}{2}m(g + a)at^{2}$，B 错误；

	CD．对集装箱进行受力分析，集装箱受到重力 $mg$ 和起重机的拉力 $F$，根据牛顿第二定律

	$F - mg = ma$

	可得起重机对集装箱的作用力 $F = m(g + a)$

	起重机的最大输出功率为 $P = Fv_{m} = m(g + a)at$，C 正确，D 错误。

	故选 AC。
}

\item
%% number: 12
%% paperName: 2025年海南高考第 12 题 4 分
%% typeId: 2
%% type: 多选题
%% chapter: 机械振动
%% point: 简谐振动
%% method: 无
%% score: 4
%% degree: 350
%% duplicateId: 0
%% body:
一绝缘的固定倾斜斜面，斜面倾角为 $30^\circ$，空间中存在沿斜面向下的匀强电场，电场强度为 $E = \frac{3mg}{2q}$。质量为 $m$ 的物块 M、N 用一根不可伸长的轻绳绕过滑轮连接，M 带正电，电荷量为 $q$，N 不带电，N 一端与弹簧连接，弹簧另一端固定在地面上，劲度系数为 $k$。初始时有外力作用使 M 静止在斜面上，轻绳恰好伸直，使 M 从静止释放，第一次到达最低点的时间为 $t$，不计一切摩擦。\xzanswer{BD}

\onepicture[w=6.20cm]{\includesvg{figs/12}}

\fourchoices[answer=BD]
{释放时 M 的加速度为 $2g$}
{M 下滑的最大速度为 $\sqrt{\frac{2mg^{2}}{k}}$}
{M 下滑的最大距离为 $\frac{3mg}{k}$}
{M 下滑的距离为 $\frac{mg}{k}$ 时，所用时间为 $\frac{t}{3}$}

%% answer: BD
%% memo:
\memoanswer{
	A．初始时，弹簧弹力等于 N 的重力，弹簧处于压缩状态，即 $kx_{1} = mg$

	可得 $x_{1} = \frac{mg}{k}$

	释放 M 时，弹簧弹力不会突变，对 M 和 N，根据牛顿第二定律

	$Eq + mg \sin 30^\circ - mg + kx_{1} = 2ma$

	可得释放时 M 的加速度为 $a = g$，A 错误；

	B．当 M、N 的加速度为零，M 的速度最大，设此时弹簧的伸长量为 $x_{2}$，根据平衡条件

	$Eq + mg \sin 30^\circ = kx_{2} + mg$

	解得 $x_{2} = \frac{mg}{k}$

	由于 $x_{1} = x_{2}$，可知弹簧弹性势能不变，M 从开始运动到速度达到最大过程，根据动能定理

	$Eq(x_{1} + x_{2}) + mg\sin 30^\circ\cdot(x_{1} + x_{2}) - mg(x_{1} + x_{2}) = \frac{1}{2}\cdot 2mv_{m}^{2} - 0$

	联立解得 M 下滑的最大速度为 $v_{m} = \sqrt{\frac{2mg^{2}}{k}}$，B 正确；

	CD．以速度最大位置为原点，斜面向上为正方向，M、N 所受的合外力与位移的关系满足

	$F_{\text{合}} = kx - (Eq + mg\sin 30^\circ - mg) = kx - mg$

	可知 M、N 做简谐运动，刚释放 M 时，加速度 $a = g$，根据简谐运动的对称性可知当 M 下滑的最大距离时，加速度大小也为 $g$，根据牛顿第二定律

	$k(x_{3} - x_{1}) + mg - Eq - mg\sin 30^\circ = 2mg$

	解得 M 下滑的最大距离为 $x_{3} = \frac{4mg}{k}$

	根据题意，M、N 做简谐运动的周期 $T = 2t$

	从释放开始计时，位移随时间变化的表达式为 $x = \frac{x_{3}}{2} \cos \frac{2\pi}{T} t' = \frac{2mg}{k} \cos \frac{\pi}{t} t'$

	当下滑距离为 $\frac{mg}{k}$ 时，代入数据有 $\frac{mg}{k} = \frac{2mg}{k} \cos \frac{\pi}{t} t'$

	可得 $t' = \frac{t}{3}$

	即 M 下滑的距离为 $\frac{mg}{k}$ 时，所用时间为 $\frac{t}{3}$，故 D 正确，C 错误。

	故选 BD。
}

\item
%% number: 13
%% paperName: 2025年海南高考第 13 题 4 分
%% typeId: 2
%% type: 多选题
%% chapter: 磁场
%% point: 带电粒子在复合场中的运动
%% method: 无
%% score: 4
%% degree: 350
%% duplicateId: 0
%% body:
某粒子分析器的部分电磁场简化模型如图，三维直角坐标系 $O-xyz$ 所在空间中 Ⅰ 区域（$0 \leq x \leq d$）存在沿 $x$ 轴正方向的匀强电场（图中未画出）和匀强磁场，磁感应强度大小为 $B_{1}$，Ⅱ 区域（$d < x < 2d$）存在沿 $z$ 轴正方向的匀强磁场，磁感应强度大小为 $B_{2}$，在 $x = 2d$ 处有一足够大的接收屏 P，原点 O 处的粒子源在 $xOz$ 平面内同时发射带正电的同种粒子甲和乙，甲粒子的速度大小为 $v_{0}$，甲和乙的速度方向与 $x$ 轴正方向夹角分别为 $30^\circ$ 和 $60^\circ$，两粒子沿 $x$ 轴方向速度分量相等。乙粒子以最短时间到达 $O_{1}(d, d, 0)$ 点进入 Ⅱ 区域后恰好到达接收屏并被吸收，不计重力及粒子间的相互作用，则\xzanswer{BD}

\onepicture[w=6.60cm]{\includesvg{figs/13}}

\fourchoices[answer=BD]
{两粒子不能同时到达接收屏 P}
{两个区域磁感应强度大小之比 $\frac{B_{2}}{B_{1}} = \frac{2}{\pi} - \frac{\sqrt{3}}{6}$}
{乙粒子通过 $O_{1}$ 点时沿 $x$ 轴方向速度分量 $v_{x} = \left(\frac{6}{\pi} + \frac{\sqrt{3}}{2}\right)v_{0}$}
{甲乙粒子在接收屏 P 上位置的 $z$ 坐标之差 $\Delta z = \frac{\pi^{2}}{12 - \sqrt{3}\pi} d$}

%% answer: BD
%% memo:
\memoanswer{
	BC．两粒子在 Ⅰ 区域运动过程，两粒子在 $x$ 轴方向做匀加速直线运动，在 $yOz$ 平面做匀速圆周运动，根据题意甲粒子和乙粒子在 $x$ 轴方向的分速度相等，均为

	$v_{\text{乙}x} = v_{\text{甲}x} = v_{0}\cos 30^\circ = \frac{\sqrt{3}}{2}v_{0}$

	甲粒子在 $z$ 轴方向的分速度 $v_{\text{甲}z} = v_{0}\sin 30^\circ = \frac{v_{0}}{2}$

	根据几何关系

	$\tan 60^\circ = \frac{v_{\text{乙}z}}{v_{\text{乙}x}}$

	可得 $v_{\text{乙}z} = \frac{3}{2}v_{0}$

	乙粒子以最短时间到达 $O_{1}(d,d,0)$，则乙在 Ⅰ 区域运动的时间为做圆周运动的周期的一半，其半径为 $r_{1} = \frac{d}{2}$

	根据洛伦兹力提供向心力

	$qB_{1}v_{\text{乙}z} = m\frac{v_{\text{乙}z}^{2}}{r_{1}}$

	联立可得 $B_{1} = \frac{3mv_{0}}{qd}$

	在 Ⅰ 区域运动的时间

	$t_{1} = \frac{\pi r_{1}}{v_{\text{乙}z}} = \frac{\pi d}{3v_{0}}$

	沿着 $x$ 正方向，根据运动学公式

	$d = \frac{v_{\text{乙}x} + v_{x}}{2}\cdot t_{1}$

	解得乙粒子通过 $O_{1}$ 点时沿 $x$ 轴方向速度分量为

	$v_{x} = \left(\frac{6}{\pi} - \frac{\sqrt{3}}{2}\right)v_{0}$

	乙粒子进入 Ⅱ 区域后，沿 $z$ 轴负方向做匀速直线运动，在 $xOy$ 平面做匀速圆周运动，根据题意乙粒子进入 Ⅱ 区域后恰好到达接收屏并被吸收，则乙粒子在 Ⅱ 区域做圆周运动的半径为 $r_{2} = d$

	根据洛伦兹力提供向心力

	$qB_{2}v_{x} = m\frac{v_{x}^{2}}{r_{2}}$

	解得

	$B_{2} = \frac{\left(\frac{6}{\pi} - \frac{\sqrt{3}}{2}\right)mv_{0}}{qd}$

	可得 $\frac{B_{2}}{B_{1}} = \frac{2}{\pi} - \frac{\sqrt{3}}{6}$，故 B 正确，C 错误；

	AD．两粒子在 Ⅰ 区域运动过程，两粒子在 $x$ 轴方向的速度分量相同，则在 Ⅰ 区域运动时间相等，根据 $T = \frac{2\pi m}{qB}$

	可知甲粒子在 Ⅰ 区域也是运动半个周期，即两粒子刚进入 Ⅱ 区域时 $z$ 轴坐标均为零，沿 $z$ 轴负方向做匀速直线运动，在 $xOy$ 平面做匀速圆周运动的情况也相同，所以运动时间相等，

	即两粒子能同时到达接收屏 P，两粒子在 Ⅱ 区域的运动时间

	$t_{2} = \frac{2\pi r_{2}}{v_{x}}\cdot\frac{1}{4}$

	甲乙粒子在接收屏 P 上位置的 $z$ 坐标之差 $\Delta z = (v_{\text{乙}z} - v_{\text{甲}z})t_{2}$

	联立解得 $\Delta z = \frac{\pi^{2}}{12 - \sqrt{3}\pi} d$，故 D 正确，A 错误。

	故选 BD。
}

\item
%% number: 14
%% paperName: 2025年海南高考第 14 题 4 分
%% typeId: 4
%% type: 实验题
%% chapter: 光学
%% point: 全反射
%% method: 无
%% score: 4
%% degree: 650
%% duplicateId: 0
%% body:
\begin{enumerate}
	\item
	小组用单摆测量当地重力加速度，装置如图~\ref{2025海南14a}所示。

	\begin{steps}
	\item
	用游标卡尺测得小球直径 $d = 20 \Umm$，刻度尺测得摆线长 $l = 79 \Ucm$，则单摆摆长 $L$ = \tkanswer{80.00} $\Ucm$（保留四位有效数字）；
	\item
	拉动小球，使摆线伸直且与竖直方向的夹角为 $\theta$（$\theta < 5^\circ$），无初速度释放小球，摆球位于\tkanswer{最低}点（选填：“最高”或“最低”）时开始计时，记录摆球做了 30 次全振动用时 $t = 54.00 \Us$，则单摆周期 $T$ = \tkanswer{1.8} $\Us$，由此测得当地重力加速度 $g$ = \tkanswer{9.88} $\Umsq$（取 $\pi^{2}$ = 10，保留三位有效数字）。
	\end{steps}
	\item
	测量某半圆形玻璃砖的折射率，操作步骤如下，装置如图~\ref{2025海南14b}所示。

	Ⅰ．在白纸上画一条直线，半圆形玻璃砖放在白纸上，玻璃砖底面直径与直线重合，描出直径两端点 A 和 B，取走玻璃砖，用刻度尺找出圆心 O 点，过 O 点作 AB 垂线 CO，放回玻璃砖，将光屏垂直 AB 贴近玻璃砖 A 点放置在玻璃砖左侧。

	Ⅱ．沿玻璃砖由 C 向 B 缓慢移动激光笔，使得入射光线平行纸面且始终沿着半径方向射向圆心 O，从玻璃砖射出的激光在 AB 下方，射在光屏上，当光屏上的光线恰好消失时，记下激光入射点 D。取走玻璃砖，过 D 点作 CO 的垂线 DE。

	（1）步骤 Ⅱ 中，当激光从 D 点入射到 O 点，光线消失的现象是光的\tkanswer{C}。

	A．色散现象\quad B．衍射现象\quad C．全反射现象

	（2）用刻度尺测得 $OB = 4.00 \Ucm$、$DE = 2.50 \Ucm$，则玻璃砖的折射率 $n$ = \tkanswer{1.6}。
\end{enumerate}

\twopicture[labelA=2025海南14a, labelB=2025海南14b, hA=3.30cm, hB=3.10cm, align=right]{figs/14a.png}{figs/14b.png}

%% answer:
\jdanswer{
\begin{enumerate}
	\item
	80.00；最低，1.8，9.88

	\item
	C；1.6
\end{enumerate}
}
%% memo:
\memoanswer{
	（1）单摆的摆长为 $L = l + \frac{d}{2} = 80.00 \Ucm$

	（2）为减小实验计时误差，需小球经过最低点时开始计时；

	单摆周期 $T = \frac{t}{n} = \frac{54}{30} \Us = 1.8 \Us$

	根据单摆周期公式 $T = 2\pi\sqrt{\frac{L}{g}}$ 可得 $g = \frac{4\pi^{2}L}{T^{2}}$，代入数值得 $g = 9.88 \Umsq$

	（1）当激光从 D 点入射到 O 点在 AB 面下方光屏上恰好消失时，此时光线在 AB 面发生全反射，故是光的全反射现象。

	故选 C。

	（2）根据前面分析可知此时入射角等于临界角，即 $\angle DOE = C$，故可得

	$\sin C = \frac{DE}{OD} = \frac{DE}{OB} = \frac{5}{8}$

	根据 $\sin C = \frac{1}{n}$

	可得玻璃砖的折射率 $n = 1.6$
}

\item
%% number: 15
%% paperName: 2025年海南高考第 15 题 6 分
%% typeId: 4
%% type: 实验题
%% chapter: 电路
%% point: 测电动势与内阻
%% method: 无
%% score: 6
%% degree: 550
%% duplicateId: 0
%% body:
某小组测量电源的电动势和内阻，可供使用的器材有：

电源 $E_{1}$（约为 $3 \UV$，内阻未知）

电压表 V（$0 \sim 3 \UV$，内阻 $R_{V}$ 约为 $3 \UkO$）

电流表 A（$0 \sim 0.6 \UA$，内阻 $R_{A}$ 约为 $1 \UO$）

定值电阻 $R_{0} = 3 \UO$

滑动变阻器 $R_{1}$（最大阻值 $50 \UO$）

滑动变阻器 $R_{2}$（最大阻值 $500 \UO$）

开关 S

导线若干

该小组设计了如图~\ref{2025海南15a}所示的实验电路：

回答下列问题：

\begin{enumerate}
	\item
	为了提高测量精度，滑动变阻器应选择\tkanswer{$R_{1}$}（选填“$R_{1}$”或“$R_{2}$”）；
	\item
	闭合开关，调节滑动变阻器，记录电压表示数 $U$ 和对应的电流表示数 $I$ 如下表，根据表格的数据在图~\ref{2025海南15b}的坐标纸上补完数据点，并绘制 $U-I$ 图像\tkanswer{（如图）}；

	\begin{center}
	\begin{tabular}{|c|c|c|c|c|c|}
	\hline
	$U/\UV$ & 1.00 & 1.30 & 1.70 & 2.00 & 2.50 \\
	\hline
	$I/\UA$ & 0.38 & 0.32 & 0.24 & 0.18 & 0.08 \\
	\hline
	\end{tabular}
	\end{center}

	\item
	由 $U-I$ 图像可求得电源电动势测量值 $E_{\text{测}}$ = \tkanswer{2.90} $\UV$，内阻测量值 $r_{\text{测}}$ = \tkanswer{2.18} $\UO$（结果均保留三位有效数字）。
	\item
	考虑电压表分流引起的误差，则 $E_{\text{测}}$ \tkanswer{小于} $E_{\text{真}}$；（填“大于”、“等于”或“小于”）；$r_{\text{测}}$ 与真实值 $r_{\text{真}}$ 之间的关系式为\tkanswer{$r_{\text{测}} = \frac{(r_{\text{真}} + R_{0})R_{\mathrm{V}}}{r_{\text{真}} + R_{0} + R_{\mathrm{V}}}$}（用 $R_{V}$，$R_{0}$，$r_{\text{测}}$，$r_{\text{真}}$ 表示）。
\end{enumerate}

\twopicture[labelA=2025海南15a, labelB=2025海南15b, wA=5.20cm, wB=6.00cm, align=right]{figs/15a.png}{figs/15b.png}

%% answer:
\jdanswer{
\begin{enumerate}
	\item
	$R_{1}$

	\item
	（如图）

	\item
	2.90，2.18

	\item
	小于，$r_{\text{测}} = \frac{(r_{\text{真}} + R_{0})R_{\mathrm{V}}}{r_{\text{真}} + R_{0} + R_{\mathrm{V}}}$
\end{enumerate}
}
%% memo:
\memoanswer{
	（1）为了提高测量精度，方便调节电路，滑动变阻器选择最大阻值较小的 $R_{1}$；

	（2）图像如下：

	\onepicture[w=6.00cm]{figs/15_an.png}

	（3）根据闭合电路欧姆定律有 $U = E_{1} - I(R_{0} + r)$

	结合图线可得 $E_{1} = 2.90 \UV$，$R_{0} + r = \frac{2.90}{0.56} \UO$

	解得 $r = 2.18 \UO$

	（4）将电压表和电源等效为新的电源，故此时根据图像得到的电动势和内阻为等效电源的电动势和内阻，等效电源的电动势为

	$E_{\text{测}} = \frac{E_{\text{真}}}{r + R_{0} + R_{V}}\cdot R_{V}$

	故 $E_{\text{测}} < E_{\text{真}}$；

	将 $R_{0}$ 等效进电源的内阻中，可得等效新电源的内阻为

	$r_{\text{测}} = \frac{(r_{\text{真}} + R_{0})R_{V}}{r_{\text{真}} + R_{0} + R_{V}}$
}

\item
%% number: 16
%% paperName: 2025年海南高考第 16 题 10 分
%% typeId: 6
%% type: 计算题
%% chapter: 气体性质
%% point: 气体多过程
%% method: 无
%% score: 10
%% degree: 650
%% duplicateId: 0
%% body:
如图，竖直放置的气缸内有一横截面积 $S = 0.01 \Umq$ 的活塞，活塞质量忽略不计，活塞与气缸壁内无摩擦且密封良好。若活塞保持静止，气缸内封闭一定质量的理想气体，气体温度 $T_{0} = 300 \UK$，体积 $V_{0} = 5 \times 10^{-3} \Umc$。设大气压强 $p_{0} = 1 \times 10^{5} \UPa$，重力加速度 $g = 10 \Umsq$。

\begin{enumerate}
	\item
	若加热气体，使活塞缓慢上升，当气体体积变为 $V_{1} = 7.5 \times 10^{-3} \Umc$，求此时气体的温度 $T_{1}$；
	\item
	若往活塞上轻放质量为 $m = 25 \Ukg$ 的重物，且活塞下降过程中气体温度 $T_{0}$ 恒定，求稳定之后的气体的体积 $V_{2}$。
\end{enumerate}

\onepicture[h=3.00cm, align=right]{figs/16.png}

%% answer:
\jdanswer{
\begin{enumerate}
	\item
	$450 \UK$

	\item
	$4 \times 10^{-3} \Umc$
\end{enumerate}
}
%% memo:
\memoanswer{
	（1）活塞缓慢上升过程中，气体做等压变化，根据盖-吕萨克定律 $\frac{V_{0}}{T_{0}} = \frac{V_{1}}{T_{1}}$

	代入数值解得 $T_{1} = 450 \UK$

	（2）设稳定后气体的压强为 $p_{2}$，根据平衡条件有 $p_{2}S = p_{0}S + mg$

	分析可知初始状态时气体压强与大气压相等为 $p_{0}$，整个过程根据玻意耳定律 $p_{0}V_{0} = p_{2}V_{2}$

	联立解得 $V_{2} = 4 \times 10^{-3} \Umc$
}

\item
%% number: 17
%% paperName: 2025年海南高考第 17 题 12 分
%% typeId: 6
%% type: 计算题
%% chapter: 运动和力
%% point: 传送带
%% method: 无
%% score: 12
%% degree: 450
%% duplicateId: 0
%% body:
某自动包装系统的部分结构简化后如图所示，足够长的传送带固定在竖直平面内，半径 $R = 0.5 \Um$，圆心角 $\theta = 53^\circ$ 的圆弧轨道与平台平滑连接，平台与顺时针匀速转动的水平传送带平滑连接，工件 A 从圆弧顶点无初速度下滑，在平台上滑入静止的空箱 B 并与其瞬间粘连成一个整体，随后一起滑上传送带，与传送带共速后进入下一道工序。已知工件 A 的质量 $m_{A} = 4 \Ukg$，空箱 B 的质量 $m_{B} = 1 \Ukg$，A、B 及粘连成的整体均可视为质点，整体与传送带间的动摩擦因数恒定，在传送带上运动的过程中因摩擦产生的热量 $Q = 2.5 \UJ$，忽略轨道及平台的摩擦，$\cos 53^\circ$ = 0.6，重力加速度 $g = 10 \Umsq$。求：

\begin{enumerate}
	\item
	工件 A 滑到圆弧最低点时受到的支持力；
	\item
	工件 A 与空箱 B 在整个碰撞过程中损失的机械能；
	\item
	传送带的速度大小。
\end{enumerate}

\onepicture[w=6.60cm, align=right]{figs/17.png}

%% answer:
\jdanswer{
\begin{enumerate}
	\item
	$72 \UN$，方向竖直向上

	\item
	$1.6 \UJ$

	\item
	$0.6 \Ums$ 或 $2.6 \Ums$
\end{enumerate}
}
%% memo:
\memoanswer{
	（1）A 从开始到滑到圆弧最低点间，根据机械能守恒

	$m_{A}g(R - R\cos 53^\circ) = \frac{1}{2}m_{A}v_{0}^{2}$

	解得 $v_{0} = 2 \Ums$

	在最低点根据牛顿第二定律

	$F_{\mathrm{N}} - m_{A}g = m_{A}\frac{v_{0}^{2}}{R}$

	解得 $F_{\mathrm{N}} = 72 \UN$，方向竖直向上；

	（2）根据题意 AB 碰后成一整体，根据动量守恒 $m_{A}v_{0} = (m_{A} + m_{B})v_{\text{共}}$

	解得 $v_{\text{共}} = 1.6 \Ums$

	故 A 与 B 整个碰撞过程中损失的机械能为

	$\Delta E = \frac{1}{2}m_{A}v_{0}^{2} - \frac{1}{2}(m_{A} + m_{B})v_{\text{共}}^{2} = 1.6 \UJ$

	（3）第一种情况，当传送带速度 $v$ 小于 $v_{\text{共}}$ 时，AB 滑上传送带后先减速后匀速运动，设 AB 与传送带间的动摩擦因数为 $\mu$，对 AB 根据牛顿第二定律 $\mu(m_{A} + m_{B})g = (m_{A} + m_{B})a$

	设经过时间 $t_{1}$ 后 AB 与传送带共速，可得 $v = v_{\text{共}} - at_{1}$

	该段时间内 AB 运动的位移为

	$x_{1} = \frac{v + v_{\text{共}}}{2}t_{1}$

	传送带运动的位移为 $x_{2} = vt_{1}$

	故可得 $Q = \mu(m_{A} + m_{B})g\cdot(x_{1} - x_{2})$

	联立解得 $v = 0.6 \Ums$，另一解大于 $v_{\text{共}}$ 舍去；

	第二种情况，当传送带速度 $v$ 大于 $v_{\text{共}}$ 时，AB 滑上传送带后先加速后匀速运动，设经过时间 $t_{2}$ 后 AB 与传送带共速，同理可得 $v = v_{\text{共}} + at_{2}$

	该段时间内 AB 运动的位移为 $x_{1}' = \frac{v + v_{\text{共}}}{2}t_{2}$

	传送带运动的位移为 $x_{2}' = vt_{2}$

	故可得 $Q = \mu(m_{A} + m_{B})g\cdot(x_{2}' - x_{1}')$

	解得 $v = 2.6 \Ums$，另一解小于 $v_{\text{共}}$ 舍去。
}

\item
%% number: 18
%% paperName: 2025年海南高考第 18 题 16 分
%% typeId: 6
%% type: 计算题
%% chapter: 电磁感应
%% point: 双杆问题
%% method: 无
%% score: 16
%% degree: 350
%% duplicateId: 0
%% body:
间距为 $L$ 的金属导轨倾斜部分光滑，水平部分粗糙且平滑相接，导轨上方接有电源和开关，倾斜导轨与水平面夹角 $\theta = 30^\circ$，处于垂直于导轨平面向上的匀强磁场中，水平导轨处于竖直向下的匀强磁场中，磁感应强度大小均为 $B$，两相同导体棒 ab、cd 与水平导轨的动摩擦因数 $\mu$ = 0.25，最大静摩擦力等于滑动摩擦力，两棒质量均 $m$，接入电路中的电阻均为 $R$，cd 棒仅在水平导轨上运动，两导体棒在运动过程中始终与导轨垂直并接触良好，且不互相碰撞，忽略金属导轨的电阻，重力加速度为 $g$。

\begin{enumerate}
	\item
	锁定水平导轨上的 cd 棒，闭合开关，ab 棒静止在倾斜导轨上，求通过 ab 棒的电流；断开开关，同时解除 cd 棒的锁定，当 ab 棒下滑距离为 $x_{0}$ 时，cd 棒开始运动，求 cd 棒从解除锁定到开始运动过程中产生的焦耳热；
	\item
	此后 ab 棒在下滑过程中电流达到稳定，求此时 ab、cd 棒的速度大小之差；
	\item
	ab 棒中电流稳定之后继续下滑，从 ab 棒到达水平导轨开始计时，$t_{1}$ 时刻 cd 棒速度为 0，加速度不为 0，此后某时刻 cd 棒的加速度为 0，速度不为 0，求从 $t_{1}$ 时刻到 cd 棒加速度为 0 的时刻 ab、cd 的路程之差。
\end{enumerate}

\onepicture[w=7.00cm, align=right]{figs/18.png}

%% answer:
\jdanswer{
\begin{enumerate}
	\item
	$\frac{mg}{2BL}$，$\frac{1}{4}mgx_{0} - \frac{m^{3}g^{2}R^{2}}{16B^{4}L^{4}}$

	\item
	$\frac{3mgR}{4B^{2}L^{2}}$

	\item
	$\Delta s = \frac{m^{2}gR^{2}}{4B^{4}L^{4}}$
\end{enumerate}
}
%% memo:
\memoanswer{
	（1）ab 棒静止在倾斜导轨上，根据平衡条件可得 $F_{\text{安}} = mg\sin 30^\circ$，$F_{\text{安}ab} = BI_{ab}L$

	解得通过 ab 棒的电流为 $I_{ab} = \frac{mg}{2BL}$

	设当 ab 棒下滑距离为 $x_{0}$ 时速度为 $v_{0}$，cd 棒开始运动时回路中的电流为 $I_{1}$，此时对 cd 棒有 $F_{\text{安}cd} = \mu mg$

	同时有 $F_{\text{安}cd} = BI_{1}L$，$I_{1} = \frac{BLv_{0}}{2R}$

	分析可知 cd 棒从解除锁定到开始运动过程中，cd 棒产生的焦耳热与 ab 棒产生的焦耳热相等，

	整个过程根据能量守恒可得

	$mgx_{0}\sin 30^\circ - \frac{1}{2}mv_{0}^{2} = 2Q_{cd}$

	联立解得 cd 棒产生的焦耳热为 $Q_{cd} = \frac{1}{4}mgx_{0} - \frac{m^{3}g^{2}R^{2}}{16B^{4}L^{4}}$

	（2）分析可知 ab 棒在下滑过程中产生的电动势与 cd 棒在向左运动的过程中产生的电动势方向相反，故当电流达到稳定时，两棒的速度差恒定，故可知此时两棒的加速度相等，由于两棒受到的安培力大小相等，对两棒有 $mg\sin 30^\circ - F_{\text{安}} = ma$，$F_{\text{安}} - \mu mg = ma$

	同时有 $F_{\text{安}} = BIL$，$I = \frac{BLv_{1} - BLv_{2}}{2R} = \frac{BL\Delta v}{2R}$

	联立解得此时 ab、cd 棒的速度大小之差为 $\Delta v = \frac{3mgR}{4B^{2}L^{2}}$

	（3）分析可知从开始到 $t_{1}$ 时刻，两棒整体所受的合外力为零，故该过程系统动量守恒，设 $t_{1}$ 时刻 ab 棒的速度为 $v_{1}'$，可知 $m\Delta v = mv_{1}'$

	解得 $v_{1}' = \Delta v = \frac{3mgR}{4B^{2}L^{2}}$

	设某时刻时，ab 棒速度为 $v_{1}''$，cd 棒速度为 $v_{cd}$，cd 棒的加速度为零，可得

	$F_{\text{安}cd2} = \mu mg$ ①

	其中 $F_{\text{安}cd2} = BI_{2}L$

	分析可知此时两导体棒产生的电动势方向相反，可得

	$I_{2} = \frac{BL(v_{1}'' - v_{cd})}{2R}$ ②

	从 $t_{1}$ 时刻到某时刻间，对两棒分别根据动量定理有 $-\left(\mu mg + B\overline{I}L\right)\Delta t = mv_{1}'' - mv_{1}'$，

	$\left(B\overline{I}L - \mu mg\right)\Delta t = mv_{cd}$

	变式可得 $\mu mg\Delta t + BLq = mv_{1}' - mv_{1}''$，$BLq - \mu mg\Delta t = mv_{cd}$

	两式相加得 $2BLq = mv_{1}' - m\left(v_{1}'' - v_{cd}\right)$ ③

	同时有 $q = \overline{I}\Delta t = \frac{\Delta\Phi}{2R}\cdot\Delta t = \frac{\Delta\Phi}{2R} = \frac{BL\Delta s}{2R}$ ④

	联立①②③④可得从 $t_{1}$ 到某时刻，ab、cd 的路程之差为 $\Delta s = \frac{m^{2}gR^{2}}{4B^{4}L^{4}}$
}

\end{enumerate}

\end{document}
